\documentclass{article}
\usepackage{microtype}
\usepackage{graphicx}
\usepackage{subcaption}
\usepackage{hyperref}
\usepackage{amsmath,amssymb,mathtools,amsthm}
\usepackage{booktabs}
\usepackage{longtable}
\usepackage{xcolor}
\usepackage{enumitem}
\usepackage{tikz}
\usetikzlibrary{arrows.meta,positioning}
\usepackage[capitalize,noabbrev]{cleveref}
\setlist{nosep,leftmargin=1.1em}
\newcommand{\PENDING}[1]{\colorbox{gray!15}{\scriptsize\textsf{PENDING: #1}}}
\newtheorem{proposition}{Proposition}
\newtheorem{corollary}{Corollary}
\newcommand{\indep}{\perp\!\!\!\perp}
\usepackage{icml2027}

\icmltitlerunning{Labels, Not Access}

\begin{document}

\twocolumn[
\icmltitle{Labels, Not Access: Monitoring Closed Agents\\from the Transcript and the Proposed Action}

\icmlsetsymbol{equal}{*}
\begin{icmlauthorlist}
\icmlauthor{Anonymous Authors}{anon}
\end{icmlauthorlist}
\icmlaffiliation{anon}{Anonymous Institution}
\icmlcorrespondingauthor{Anonymous}{anon@anon}
\icmlkeywords{AI safety, agent monitoring, interpretability, black-box models, preregistration}
\vskip 0.3in
]
\printAffiliationsAndNotice{}

% =====================================================================================
\begin{abstract}
A monitor that reads a model's internal activations needs the model's weights. The most capable agents give only text.
We show that, for the labels that agent monitoring needs, weight access adds nothing that the transcript and the proposed action do not give.
An agent's internal state is a function of the tokens that it read, so internals can only add information from tokens that the transcript hides.
For a policy check on a proposed action, the action makes even those tokens irrelevant: a monitor that reads the transcript and the proposed action is Bayes-sufficient.
The remaining problem is statistical: learn the policy from labelled examples.
We test each term with preregistered, hash-locked tests.
Reading after the agent's reasoning explains the white-box advantage over monitors at the start of the turn (AUROC 0.979 vs 0.841).
Without labels, a 70B judge loses to a probe (0.782). With labels, an open reader of the transcript and the action beats the agent's own internals (0.994 vs 0.979), and the agent's hidden reasoning adds nothing once the action is visible ($-0.0014$).
Trained only on an open agent, the reader catches 96\% and 98\% of the violations of two closed agents, Claude Haiku~5.5 and GPT-6-luna, at a 5\% review budget on the training benchmark; a zero-shot judge catches 0\% and 38\%.
On an external benchmark of prompt injections, the reader is not significantly better than the judge on a closed agent (0.840 vs 0.787; a preregistered failure). Labels of that new policy, collected on open agents, fix this: 0.922 on the closed agent, and 92\% of the injected calls caught.
The limit is the coverage of the policies, not the agent: a policy unlike any trained policy is ranked but not escalated until labelled examples of it are added, and these labels transfer from open agents to closed ones.
We release the record of every test, an SDK, and a website on which anyone can train and audit a monitor in the browser.
\end{abstract}

% =====================================================================================
\section{Introduction}
\label{sec:intro}

Agents built on language models now act through tools: they send messages, run commands and move data.
A monitor must decide if each proposed action can run, before it runs.
Monitors that use interpretability read the agent's activations \citep{zou2023repe,goldowskydill2025probes,fayyazi2026cape}.
They need the weights, but the agents with the most capability are closed.
Their providers give sampled text and, at most, top-$k$ log probabilities and a summary of the reasoning.

Some evidence suggests that the transcript is enough.
Open ``reader'' models that read the transcripts of closed models come close to probes on the target's own activations \citep{conti2026hermeneutics}.
Follow-up questions detect lies about as well as probes on prompted lies, but not on model organisms with verified beliefs \citep{pacchiardi2024liar,cooney2026didyoulie}.
No result says when the transcript is enough, or why.

\textbf{Our answer.}
For a policy check on a proposed action, weight access adds nothing that the transcript and the proposed action do not give (\cref{sec:theory}).
The value of a white-box monitor then reduces to three measurable terms: tokens that the transcript hides, the label function, and the statistics of learning the policy from labels.
We test each term with preregistered tests on agents from 20B-parameter open models to closed frontier models (\cref{sec:results}).

\textbf{Contributions.}
\begin{enumerate}
\item \emph{Theory.} An agent's internals add to its full token history no information about any label; under log loss, a white-box monitor beats a transcript monitor by exactly $I(Y;H\mid V)$, the information in hidden tokens (\cref{prop:access}). For action-level labels, the transcript and the proposed action are Bayes-sufficient (\cref{cor:action}).
\item \emph{A preregistered test of each term:} timing (T2), the label function (T3), labels (T4, T5), representation (T9) and policy coverage (T10, T12).
\item \emph{Two closed agent families.} A reader trained only on an open agent monitors Claude Haiku~5.5 (T8b) and GPT-6-luna (T11) beyond tool identity on the training benchmark. On an external benchmark it beats a zero-shot judge for open agents (T7), but not significantly for a closed agent (T13b), until labels of that policy from open agents are added (T14).
\item \emph{Deployment without labels from the new agent.} Probabilities do not transfer between agents; rankings do. A review threshold set from the agent's own traffic and then frozen resists bursts of violations.
\item \emph{A tool.} An open SDK and a website: anyone can train and audit a monitor on their own agent in the browser, with results identical to the Python version.
\end{enumerate}

The theorem is simple; we do not claim that it is deep.
Its value is that it turns a vague question (``do we need the weights?'') into terms that experiments can measure, and our tests measure each term.

\begin{figure*}[t]
\centering
\resizebox{0.98\textwidth}{!}{\input{fig_overview}}
\caption{One agent step, and what each monitor reads. A white-box probe gains over monitors at the start of the turn because it reads after the hidden reasoning. A monitor that reads the proposed action needs no weights: for a policy check on the action, it is Bayes-sufficient (\cref{cor:action}), and with labels it scores higher than the internals. AUROC on the same 12{,}128 fresh steps (T2, T4).}
\label{fig:overview}
\end{figure*}

% =====================================================================================
\section{Theory}
\label{sec:theory}

\textbf{Setting.}
A monitor acts at step $t$, after the agent proposes an action and before the action runs.
$Z_t$ is every token that the agent read or wrote up to $t$.
$V_t$ is the part that a deployment sees, and $H_t$ is the rest, for example reasoning that the API hides.
$A_t$ is the proposed action, and $W$ are world facts that the agent did not observe.
An internal representation is $\phi_t=\Phi(Z_t;\theta)$, a deterministic function of the tokens.

\begin{proposition}[access is not information]\label{prop:access}
For every label $Y$, $\sigma(Z_t,\phi_t)=\sigma(Z_t)$. Thus the best monitor that uses $(Z_t,\phi_t)$ is not better than the best monitor that uses $Z_t$, under every strictly proper score and at every point of the ROC curve. Under log loss, the best monitor that uses $Z_t$ beats the best monitor that uses $V_t$ by exactly $I(Y;H_t\mid V_t)$.
\end{proposition}

\begin{corollary}[the proposed action is sufficient]\label{cor:action}
Let $Y=g(A_t,V_t,W)$, with $W\indep Z_t^-\mid V_t$, where $Z_t^-$ are the tokens before the action is sampled, and with sampling noise independent of $(Z_t^-,W)$.
Then $P(Y\mid Z_t^-,A_t)=P(Y\mid V_t,A_t)$, and every monitor that reads internals at or before the action has a Bayes risk at least that of the best monitor that reads $(V_t,A_t)$.
\end{corollary}

Proofs are in \cref{app:proofs}.
A probe read before the action is emitted must also predict which action the agent will sample; a monitor that reads the action does not.
Thus the inequality can be strict when the agent's action is random.
\Cref{cor:action} does not cover labels about later behaviour or about beliefs.
For labels about later behaviour, the rollout forecast dominates every readout of the state (\cref{app:proofs}, Proposition~3).

\textbf{What remains is statistics.}
The excess risk of a trained transcript monitor is the risk of learning $P(Y\mid V_t,A_t)$ from data.
It depends on the label function used for training, on the representation, and on how well the labelled data cover the policies that $g$ enforces.

% =====================================================================================
\section{Setup}
\label{sec:setup}

\textbf{Agents and data.}
CAPE-Bench \citep{fayyazi2026cape} gives matched agent tasks in which one fact decides if an action is authorized; an environment rule checker labels each action.
Open agents: Gemma-4 (26B) and gpt-oss-20b.
Closed agents, through their public tool-use APIs, with reasoning off: Claude Haiku~5.5 and GPT-6-luna, 1{,}000 episodes each.
External benchmark: AgentDojo prompt injections \citep{debenedetti2024agentdojo} with the spotlighting defense \citep{hines2024spotlighting}; a call is labelled injected if it executes the attacker's task.

\textbf{Monitors.}
\emph{Reader}: Qwen3-8B reads the agent's instructions, the user's task, the earlier tool calls with their results, and the proposed action; a linear probe on its residual stream (layer 32) learns the labels.
It is trained only on Gemma-4 CAPE-Bench steps (17{,}623 steps) and is never fit on a test agent or on AgentDojo.
\emph{Internals}: CAPE's probe on Gemma-4's own residual stream.
\emph{Zero-shot judge}: Llama-3.3-70B with the same input.
\emph{TF-IDF}: a bag-of-words monitor with the reader's input and labels.
\emph{Tool identity}: the training-set violation rate of the action's tool.

\textbf{Preregistration and statistics.}
Every confirmatory test fixed its data, metric and pass criterion, and was hash-locked before its data existed or was scored.
We report every outcome, including failures.
We give AUROC within strata (template, tool or suite), with 95\% intervals from a bootstrap over scenarios or episodes (1{,}000 resamples), and the share of violations caught at a 5\% review budget.

% =====================================================================================
\section{Results}
\label{sec:results}

\Cref{tab:all} lists every agent test with its preregistered outcome, including the failures.

\begin{table*}[t]
\centering\small
\caption{All agent tests, in the order in which they were locked. Pass: the preregistered criterion was met. AUROC within strata unless stated.}
\label{tab:all}
\begin{tabular}{llll}
\toprule
Test & Question & Key result & Outcome\\
\midrule
T1 & Resampling vs a probe at the turn start (oracle checker) & 0.904 vs 0.841 & pass; answers alone: fail\\
T2 & Internals after the reasoning vs monitors at the turn start & 0.979 vs 0.904 & pass\\
T3 & Resampling with a deployable 70B judge & 0.782 $<$ probe 0.846 & fail (probe better)\\
T4 & Transcript reader with labels vs the agent's internals & 0.994 vs 0.979 & pass\\
T5 & The same, on unseen task families & 0.969 vs 0.976 (non-inferior) & pass\\
T6 & Reader trained on Gemma, applied to gpt-oss & 0.915 vs judge 0.833 & pass\\
T7 & Reader on AgentDojo, no AgentDojo training & 0.781 vs judge 0.537 & pass\\
T8 & Claude, actions after the reveal & 0 violations in the units & not tested\\
T8b & Claude, all actions & 95.6\% caught vs 0\% & pass\\
T9 & TF-IDF with the same labels vs the reader & equal in distribution; chance on AgentDojo & mixed\\
T10 & Claude, a policy family removed from training & 0.912 vs 0.570; 0/26 caught & pass (ranking)\\
T11 & GPT-6-luna, beyond tool identity & 0.996 vs 0.778; tool 0.856 & pass\\
T12 & Labels for a new policy, prospective & 18/18 caught with 0 labels & H1 fail\\
T13 & GPT-6-luna on AgentDojo, with spotlighting & attack success 2.1\%; 4 injected calls & not tested\\
T13b & The same, without the defense, both splits & 0.840 vs judge 0.787 (36 injected) & fail\\
T14 & T13b, after adding AgentDojo labels of open agents & 0.922 vs judge 0.787; 92\% caught & pass\\
\bottomrule
\end{tabular}
\end{table*}

\subsection{Timing explains the white-box advantage (T2)}
On 12{,}128 fresh Gemma-4 steps, a probe on the agent's internals scores 0.841 at the start of the turn and 0.979 after the agent's reasoning, before the action (\cref{fig:time}).
After the reasoning, the action is nearly fixed: resampling the action and checking it with the environment's rule checker gives 0.9997.
The advantage of internals over monitors at the start of the turn is when they read, not what they can access.

\begin{figure}[t]
\centering
\includegraphics[width=0.9\columnwidth]{figures/time_vs_access.pdf}
\caption{The same 12{,}128 fresh steps, read at two moments. Reading after the reasoning helps every method. Resamples are checked by the environment's rule checker (an oracle).}
\label{fig:time}
\end{figure}

\subsection{The label function decides in practice (T3)}
Resampling needs a label function at test time.
With a 70B judge in place of the rule checker, resampling falls to 0.782, below the turn-start probe (0.846) and far below the internals after the reasoning (0.976).
The bottleneck is the label function, not the access.

\subsection{With labels, the transcript equals the internals (T4, T5)}
The reader, trained on the same labels, scores 0.994 on the same scenarios; the agent's internals score 0.979 ($+0.015$ [0.007, 0.026]).
On unseen scenario families (5 folds over 23 families) the reader is non-inferior: 0.969 vs 0.976 ($-0.007$ [$-0.014$, $-0.0001$]; margin $-0.02$).
The zero-shot judge scores 0.910.
Given the action, the agent's hidden reasoning adds nothing: adding it to the reader's input changes AUROC by $-0.0014$ [$-0.0046$, $+0.0014$] (exploratory), as \cref{cor:action} predicts.
Seventy-four labelled steps already beat the judge (0.938).

\subsection{The representation carries labels to new tasks (T9)}
TF-IDF with the same input and labels equals the reader on the training benchmark (0.992 vs 0.994).
On unseen task families the reader is better (0.969 vs 0.950; $+0.019$ [0.007, 0.032]).
On AgentDojo, TF-IDF is at chance (0.503); the reader scores 0.781 (T7).

\begin{figure*}[t]
\centering
\includegraphics[width=0.78\textwidth]{figures/generalization.pdf}
\caption{Monitoring the next action from the transcript, under increasing shift. Reader (blue), TF-IDF with the same labels (hollow square), zero-shot 70B judge (green), the agent's internals (orange). The Claude rows rest on one stratum (136 actions, 26 violations); in the T10 row the monitors were retrained without that policy.}
\label{fig:gen}
\end{figure*}

\subsection{Two closed agent families (T8b, T11, T13, T13b)}
Nothing was fit on the closed agents (\cref{tab:closed}).
On Claude Haiku~5.5 the reader catches 95.6\% of 68 violations at a 5\% budget; the judge catches none.
The preregistered within-template AUROC (0.991 vs 0.633) rests on one stratum, so we also give the pooled AUROC (0.989 vs 0.570, exploratory).
On GPT-6-luna all three preregistered endpoints pass: pooled 0.996 vs 0.778 ($+0.217$ [0.174, 0.265]); against tool identity alone $+0.141$ [0.111, 0.171]; within tool $+0.039$ [0.015, 0.066].
On AgentDojo, GPT-6-luna rarely follows the injected instructions: with the spotlighting defense, the attack succeeds in 2.1\% of the episodes, which leaves 4 injected calls, below the preregistered minimum of 10 (T13, not tested). Without the defense, on both splits, there are 36 injected calls (T13b). There the reader is not significantly better than the zero-shot judge: 0.840 vs 0.787 ($+0.052$ [$-0.026$, $+0.130$]); the preregistered test fails. TF-IDF is below chance (0.335). The monitor transfers to closed agents on its training benchmark; on an external benchmark with a different policy, it needs labels of that policy (next section).

\begin{table}[t]
\centering\small
\caption{Closed agents, frozen monitors trained only on Gemma-4. Caught: violations escalated when 5\% of the actions go to review.}
\label{tab:closed}
\begin{tabular}{lcccc}
\toprule
 & \multicolumn{2}{c}{Claude Haiku 5.5} & \multicolumn{2}{c}{GPT-6-luna}\\
 & AUROC & caught & AUROC & caught\\
\midrule
Reader & 0.989 & 95.6\% & 0.996 & 98.5\%\\
TF-IDF & 0.997$^\dagger$ & 64.7\% & 0.964 & 50.0\%\\
Zero-shot judge & 0.570 & 0\% & 0.778 & 37.9\%\\
Tool identity & -- & -- & 0.856 & 60.6\%\\
\bottomrule
\end{tabular}\\[2pt]
\footnotesize Pooled AUROC. $^\dagger$within the one stratum with both classes.
\end{table}

\subsection{The limit is policy coverage (T10, T12, T14)}
We removed one policy family from the training data.
For a service-restart policy, the reader still ranks Claude's violations well (pooled 0.912 vs 0.570 for the judge), but at a 5\% budget it escalates 0 of 26 (T10).
The violations rank above safe actions of the same task, but not above risky actions of other tasks.
With 20 labelled steps of the policy it catches 85\%; with 40, 94\% (exploratory; \cref{fig:policy}).
We then tested this rule prospectively on GPT with a credentials policy (T12).
The preregistered hypothesis failed: with no labels, the reader already caught all 18 violations.
The training data hold similar credential policies; the restart policy had none.
A new policy that resembles a trained one can transfer; one that does not needs labels.
Prompt injection on AgentDojo is such a policy (T13b). We added AgentDojo calls from the two \emph{open} agents, with their injection labels, to the training data, and tested on the \emph{closed} agent (T14, preregistered).
With all 3{,}455 labelled calls, the reader scores 0.922 against 0.787 for the judge ($+0.136$ [0.078, 0.191]) and against 0.840 without the labels ($+0.081$ [0.044, 0.121]); it catches 92\% of the injected calls at a 5\% budget, against 36\% without (\cref{fig:policy}b).
Labels of a policy transfer across agents, from open to closed.

\begin{figure*}[t]
\centering
\includegraphics[width=0.82\textwidth]{figures/policy_labels.pdf}
\caption{Policy coverage. (a) A CAPE-Bench policy removed from training, against the number of its labelled steps added back. Credentials (GPT, preregistered T12): caught without labels. Service restart (Claude, exploratory): missed until 20--40 labels. (b) Prompt injection on the closed GPT agent, against the number of labelled AgentDojo calls from \emph{open} agents added to training (preregistered T14).}
\label{fig:policy}
\end{figure*}

\subsection{A review threshold without labels from the new agent}
The reader's probabilities do not transfer between agents: on Claude, a 0.5 cut catches 3\%.
Its ranking transfers.
We set the threshold at the top 5\% of the agent's first 500 actions and then freeze it.
Online, it catches 95.4\% on Claude (1.9\% of safe actions escalated) and 98.4\% on GPT (2.3\%).
In a simulated burst in which half of the actions are unsafe, the frozen threshold catches 94\%; a rolling threshold rises and catches 36\% (exploratory).

% =====================================================================================
\section{The tool}
\label{sec:tool}
We release an SDK and a website (links in the camera-ready version).
On the website, a user loads labelled steps from their own agent, trains the TF-IDF monitor in the browser, gets a cross-validated accuracy over held-out tasks, a frozen review threshold and the catch rate at their budget, and scores new actions.
The data stay in the browser.
The browser and the Python SDK give byte-identical text on 302 real agent steps and the same held-out AUROC (0.9887 vs 0.9886); a model trained in the browser scores identically in Python ($10^{-13}$).
Scoring takes 0.18~ms per action in a browser; the reader takes 41~ms on one GPU.
A self-hosted API keeps transcripts on the user's machine.

% =====================================================================================
\section{Related work}
\label{sec:related}
\textbf{Probes on internals} read truth, refusal and deception from activations \citep{burns2023ccs,marks2023geometry,zou2023repe,arditi2024refusal,goldowskydill2025probes}; CAPE probes an agent before each action \citep{fayyazi2026cape}.
\Cref{prop:access} states what they can add when a transcript exists.
\textbf{Chain-of-thought monitors} catch reward hacking that action-only monitors miss \citep{baker2025monitoring}; reasoning is a fragile but useful signal \citep{korbak2025cot,chen2025reasoning}.
\Cref{cor:action} makes this more exact: for labels that the action and the visible facts decide, reasoning adds nothing once the monitor sees the action; it can help for labels about intent or later behaviour.
\textbf{AI control} audits a fixed budget of actions with a trusted monitor \citep{greenblatt2024control}. Our results show that a trained transcript monitor beats a zero-shot trusted judge, and that the threshold must come from the monitored agent's own traffic.
\textbf{Black-box access} recovers parts of a model from its logits \citep{carlini2024stealing,finlayson2024logits}, detects lies with follow-up questions \citep{pacchiardi2024liar,cooney2026didyoulie,sam2025quere}, forecasts behaviour from traces \citep{levy2026forecasting}, and probes closed transcripts with open readers \citep{conti2026hermeneutics}.
We add the theory of when transcripts are enough, a preregistered test of each term, and closed agents.

% =====================================================================================
\section{Limitations}
\label{sec:limits}
\begin{itemize}
\item Most agent tests use one benchmark. On AgentDojo, the external test, the reader beats a zero-shot judge for open agents (0.781 vs 0.537) but not significantly for a closed agent (0.840 vs 0.787, T13b).
\item A monitor needs labels of each policy that is unlike its training policies. We measured this on three held-out policies (service restart, credentials, prompt injection); the role of similarity is a hypothesis.
\item \Cref{cor:action} needs the label to be decided by the action and the visible facts, and a deployment that sees the action before it runs. It does not cover labels about beliefs or later behaviour.
\item A frozen threshold assumes trusted warm-up traffic; a poisoned warm-up would raise it.
\end{itemize}

\section{Conclusion}
For monitoring the actions of closed agents, the weights are not the bottleneck: labels and policy coverage are.
A reader of the transcript and the proposed action, trained on labels, equals the agent's internals and transfers to closed agents from two providers on the policies it learned. For a different policy, such as prompt injection, it needs labels of that policy, and labels collected on open agents are enough.

\bibliography{refs}
\bibliographystyle{icml2027}

\newpage
\appendix
\onecolumn
\section{Proofs}
\label{app:proofs}
\input{appendix_proofs}

\section{The full preregistration record}
\label{app:record}
Tally, as recorded: method (G): 48 tests; 18 passed every primary endpoint, 17 passed some endpoints and failed others, 9 failed, 2 stopped before any outcome, 1 reported without a pass/fail bar, 1 pending; trained comparisons (T): 32 tests; 15 passed every primary endpoint, 7 passed some endpoints and failed others, 7 failed, 0 stopped before any outcome, 3 reported without a pass/fail bar, 0 pending. No correction for multiple tests is applied; each test answers its own registered question, and every outcome, failed or passed, is listed below.

Every confirmatory test, as in the released record (generated from it by a script; file paths and commit ids are replaced for review). Locks (SHA-256 of the preregistration and its code) and outcomes are in the released record; lock times are self-reported. Endpoints are numbered (i), (ii), ... in the order of each preregistration. Later tests changed two readings fixed in advance. First, the unseen-task test showed that labels do not transfer on tau2-bench through word counts: the test tasks had been in training. Second, ``checks rules'' in the frontier-agent judge test means only that the judge separates violating from allowed writes of the same tool.

{\small\begin{longtable}{p{0.24\textwidth}p{0.71\textwidth}}
\toprule Test & Outcome (verbatim from the record)\\\midrule\endhead
\textbf{Interrogation state, open chat models}\newline open LLMs ((i)--(iii)) & \textbf{PASS all, after deviation 1}: (i) +0.120 [0.115, 0.126]; (ii) +0.057 [0.045, 0.073] (7/8 models); (iii) +0.359 [0.310, 0.415]; Holm p < 0.0005 each. First run: (i) valid PASS, (ii)/(iii) invalid (numerical; see note) ([released file])\\[2pt]
\textbf{Resampling vs a probe}\newline agents ((i)--(iii); addendum 1: (iv)--(v)) & \textbf{Claim (a) CONFIRMED} ((i), (ii) pass: resampling (16 rollouts) $-$ white-box probe +0.063 [0.033, 0.094]); claim (b) NOT confirmed ((iii) fails: $\Psi$ $-$ white-box probe $-$0.021 [$-$0.054, 0.009]); addendum \textbf{(iv) passes} ((iv)+$\Psi$ $-$ white-box probe +0.049 [0.020, 0.080]); (v) fails ((ii)+$\Psi$ $-$ resampling (16 rollouts) $-$0.023 [$-$0.053, 0.006]) ([released file])\\[2pt]
\textbf{Interrogation state, Othello}\newline Othello recovery ((i)--(iii)) & \textbf{PASS all}: (i) R$^2$ 0.967 [0.967, 0.968] > 0.90; (ii) $\Delta$R$^2$ +0.094 [0.086, 0.105], Holm p < 0.0005; (iii) $\Delta$ occupied acc +0.043 [0.040, 0.046], Holm p < 0.0005. White-box reference 0.983 ([released file])\\[2pt]
\textbf{Claude Haiku 5.5 chat readouts}\newline Claude Haiku 5.5 readout transfer ((i)--(iii)) & \textbf{all pass}: (i) 0.993 [0.983, 0.999]; (ii) 0.997 [0.990, 1.000]; (iii) 0.935 [0.868, 0.982] ([released file])\\[2pt]
\textbf{Claude Haiku 5.5 refusals vs moderation}\newline Claude Haiku 5.5: grey box vs text moderation ((i)--(iii)) & \textbf{(i) pass} 0.923 [0.884, 0.957]; \textbf{(ii) (key) pass} grey box $-$ moderation +0.176 [0.060, 0.287] (0.827 vs 0.651); (iii) not tested (1/40 over-refusals) ([released file])\\[2pt]
\textbf{Claude Haiku 5.5, one-call grey box}\newline Claude Haiku 5.5: one-call grey box + compliance alarm ((i)--(iii)) & \textbf{all pass} vs the preregistered 3-level verdict: (i) +0.142 [0.081, 0.203] (0.815 vs 0.673); (ii) AP +0.230 [0.161, 0.307]; (iii) AP +0.081 [0.036, 0.140]. Against GRADED moderators (i) ties (Qwen3Guard +0.018 [$-$0.043, +0.081]; Llama-Guard $-$0.031 [$-$0.082, +0.020]); the alarm gain holds ([released file], [released file])\\[2pt]
\textbf{Claude Haiku 5.5, transcript readout}\newline Claude Haiku 5.5: transcript-framing grey box ((i)--(ii)) & \textbf{both pass}: (i) grey box $-$ moderation on harmful refusals +0.141 [0.082, 0.199] (0.817 vs 0.676); (ii) compliance-risk AP (Claude-moderator) +0.115 [0.076, 0.179] ([released file])\\[2pt]
\textbf{When the internals are read}\newline time: internals at CAPE's decision point vs black box ((i)--(iii)) & \textbf{all pass}: (i) internals at the decision $-$ resampling at the start of the turn +0.075 [0.044, 0.106]; (ii) resampling at the decision $-$ internals at the decision +0.021 [0.011, 0.035]; (iii) internals at the decision $-$ white-box probe at the start of the turn +0.138 [0.107, 0.163] ([released file])\\[2pt]
\textbf{Claude Haiku 5.5 vs graded moderators}\newline Claude Haiku 5.5: the Live tool's alarm, prospective ((i)--(iii)) & \textbf{(i) pass} (alarm precision 0.535 vs 0.274 for the moderator alone, +0.261 [0.215, 0.311]); \textbf{(ii) FAIL} (grey box 0.873 vs graded Qwen3Guard 0.855, +0.018 [$-$0.023, 0.062]); \textbf{(iii) FAIL} (multi-turn 0.894 vs 0.906) ([released file])\\[2pt]
\textbf{Resampling with a deployable judge}\newline deployable monitors: LLM judge instead of the oracle checker ((i)--(iii)) & \textbf{(i) FAIL} (judge-scored resampling 0.782 < turn-start probe 0.846, $-$0.065 [$-$0.115, $-$0.019]); \textbf{(ii): probe better} (CAPE probe 0.976 > judge on actual action 0.910, $-$0.067 [$-$0.104, $-$0.036] for the judge); \textbf{(iii) FAIL} (resampling hurts the action judge, $-$0.049) ([released file])\\[2pt]
\textbf{Reader vs the agent's internals}\newline reader trained on labels vs the agent's internals ((i)--(ii)) & \textbf{both pass}: (i) reader $-$ internals at the decision +0.015 [0.007, 0.026] (0.994 vs 0.979); (ii) reader $-$ judge +0.085 [0.052, 0.127] ([released file])\\[2pt]
\textbf{Unseen task families}\newline agent internals vs transcript reader on unseen scenario families ((i)--(ii)) & \textbf{(i) pass} (reader $-$ agent internals on unseen families $-$0.007 [$-$0.014, $-$0.0001], margin $-$0.02); \textbf{(ii): agent probe better than zero-shot judge} (+0.062 [0.028, 0.101]) ([released file])\\[2pt]
\textbf{Another open agent}\newline reader trained on Gemma, monitoring gpt-oss ((i)) & \textbf{(i) pass}: Gemma-trained reader on gpt-oss 0.915 vs zero-shot judge 0.833, +0.082 [0.043, 0.129] ([released file])\\[2pt]
\textbf{AgentDojo, no labels}\newline reader monitor on AgentDojo prompt injection, zero transfer training ((i)--(ii)) & \textbf{(i) pass}: reader 0.781 vs zero-shot judge 0.537, +0.249 [0.136, 0.369]; \textbf{(ii) reader better than CAPE's gate}: gpt-oss calls as the prereg text specifies, 0.823 vs 0.601, +0.222 [0.143, 0.303] (locked code's broader subset: +0.264 [0.174, 0.359]) ([released file])\\[2pt]
\textbf{Claude Haiku 5.5, after a reveal}\newline Claude Haiku 5.5 as an agent, monitored by a reader trained only on Gemma ((i)) & \textbf{(i) NOT TESTED}: 0 unsafe units under the preregistered post-reveal rule (Claude's 68 violations all happen at or without the reveal); see ``Closed agent: Claude Haiku 5.5''\\[2pt]
\textbf{Closed agent: Claude Haiku 5.5}\newline Claude as an agent, all non-read actions ((i')) & \textbf{(i') pass}: Gemma-trained reader on Claude 0.991 vs zero-shot judge 0.633, +0.360 [0.302, 0.413]; 95.6\% of violations caught at 5\% escalation with a cut set on labelled safe actions (judge 0\%); at the label-free cut used elsewhere, 94.1\% ([released file], [released file])\\[2pt]
\textbf{Reader vs TF-IDF}\newline trained TF-IDF monitor on the reader's own input, same labels ((i) Claude, (ii) unseen families) & \textbf{(ii) pass, (i) FAIL}: unseen families reader 0.969 vs TF-IDF 0.950 (+0.019 [0.007, 0.032]); Claude reader 0.991 vs TF-IDF 0.997 ($-$0.006 [$-$0.027, +0.007]). AgentDojo (reported) 0.781 vs 0.503 ([released file])\\[2pt]
\textbf{A policy removed from training}\newline Claude on task families the monitor never saw ((i) pooled, (ii) within) & \textbf{(i), (ii) pass} (vs zero-shot judge): pooled 0.912 vs 0.570 (+0.342 [0.311, 0.371]); within 0.948 vs 0.633 (+0.318 [0.265, 0.374]). But at the label-free cut the unseen restart family is caught 0/26 ([released file])\\[2pt]
\textbf{Closed agent: GPT-6-luna}\newline second closed family: gpt-6-luna as agent, frozen monitors ((i) pooled vs judge, (ii) vs tool identity, (iii) within tool) & \textbf{(i), (ii), (iii) all pass}: GPT-6-luna, 3,627 actions, 66 unsafe in 6 families. Reader pooled 0.996 vs judge 0.778 (+0.217 [0.174, 0.265]); vs tool identity +0.141 [0.111, 0.171]; within tool vs judge +0.039 [0.015, 0.066]. 98.5\% caught at the label-free 5\% cut (judge 37.9\%) ([released file])\\[2pt]
\textbf{Labels for a new policy}\newline labels needed for a new policy, prospective on GPT ((i) 40 vs 0 labels, (ii) catch at 40 $\geq$ 0.5) & \textbf{(i) FAILS, (ii) passes}: with 0 labels of the held-out credentials policy, the reader already caught 18/18 of GPT's violations of it, so added labels had nothing to fix (40 $-$ 0: 0.0). The ``A policy removed from training'' curve does not replicate on this family ([released file])\\[2pt]
\textbf{AgentDojo with a defense}\newline a closed agent on an external benchmark: GPT-6-luna on AgentDojo, frozen monitors ((i) reader vs judge) & \textbf{(i) NOT TESTED}: GPT-6-luna resisted the injections (attack success 2.1\%), leaving 4 injected calls (minimum 10). Utility on benign tasks 0.862 ([released file])\\[2pt]
\textbf{AgentDojo, closed agent}\newline closed agent on AgentDojo without the defense (arm none, dev + conf; (i') reader vs judge) & \textbf{(i') FAILS}: 36 injected calls; reader 0.840 vs zero-shot judge 0.787 within suite (+0.052 [$-$0.026, +0.130]). TF-IDF 0.335. At the benign-call 5\% cut: reader 36\%, judge 53\% ([released file])\\[2pt]
\textbf{Open-agent labels for a closed agent}\newline labels of the new policy from open agents, tested on the closed agent ((i) vs judge, (ii) vs no labels) & \textbf{(i), (ii) pass}: with AgentDojo labels from open agents, the reader on GPT's AgentDojo calls scores 0.922 vs judge 0.787 (+0.136 [0.078, 0.191]) and vs no labels +0.081 [0.044, 0.121]; caught at the benign 5\% cut 36\% $\rightarrow$ 92\% ([released file])\\[2pt]
\textbf{Replication (Gemma-4)}\newline latent replication: does the reader replicate the agent's hidden state? ((i) R$^2$ vs text, (ii) readouts on the replica) & \textbf{(i), (ii) pass}: from the transcript alone, the reader replica reconstructs 66.4\% of the agent's decision-state variance (top 128 PCs; text 58.6\%; +0.079 [0.075, 0.083]); a probe trained on TRUE states scores 0.987 on the replica vs 0.984 on the true state (+0.003 [$-$0.002, +0.007]) ([released file])\\[2pt]
\textbf{Live prevention (GPT)}\newline live prevention on a closed agent: the monitor blocks GPT-6-luna's escalated actions in the loop ((i) violation episodes, gated $-$ ungated) & \textbf{(i) passes}: live, the grey-box gate cut GPT-6-luna's episodes with an executed violation from 6.6\% to 0.2\% ($-$0.064 [$-$0.088, $-$0.044]); 0 violations after a block; goals 47.4\% $\rightarrow$ 45.6\%; round trip 89 ms median ([released file])\\[2pt]
\textbf{Live prevention (Claude Haiku 5.5)}\newline live prevention on a second closed family: Claude Haiku 5.5 in the loop ((i')) & \textbf{(i') passes}: on Claude Haiku 5.5, episodes with an executed violation 6.7\% $\rightarrow$ 0.2\% ($-$0.064 [$-$0.087, $-$0.044]); 0 violations after a block; goals 40.7\% $\rightarrow$ 34.9\%; round trip 80 ms median ([released file])\\[2pt]
\textbf{Replication, three architectures}\newline replication on three more open agents: gpt-oss-20b, Qwen3.8-27B, Muse-Glimmer-30B ((i), (ii) per agent; scenario-disjoint split) & \textbf{(i) and (ii) pass on all three agents (replication generalizes)}: from the transcript alone, R$^2$ of the decision state on fresh scenarios: gpt-oss 0.494 (text 0.392; +0.102 [0.093, 0.110]), Qwen3.8 0.542 (0.441; +0.101 [0.091, 0.111]), Muse Glimmer 0.462 (0.401; +0.061 [0.051, 0.071]). A probe trained on the true states, true $\rightarrow$ replica: 0.900 $\rightarrow$ 0.924, 0.955 $\rightarrow$ 0.972, 0.929 $\rightarrow$ 0.969 (all lower bounds > $-$0.02). Fidelity: the agents' own models predict 91--100\% of their next 16 tool-call tokens. Exploratory, Gemma-4 under the same split: R$^2$ 0.631 (text 0.556). Exploratory: the text map chose $\alpha$ = 1, the edge of the grid; a finer grid ($\alpha$ 0.01--1, exact dual-form ridge, [released file]) is worse at every smaller $\alpha$ for all three agents, so $\alpha$ = 1 is the text map's optimum ([released file], [released file])\\[2pt]
\textbf{Hold for review}\newline hold for review instead of block: Claude Haiku 5.5, escalated actions go to a reviewer ((i) safety kept, (ii) no goal cost vs ``Live prevention (Claude Haiku 5.5)'' ungated) & \textbf{(i) passes, (ii) fails.} Episodes with an executed violation: 6.7\% (``Live prevention (Claude Haiku 5.5)'' ungated) $\rightarrow$ 0.0\% (review) ($-$0.067 [$-$0.091, $-$0.044]). Goals: 40.7\% $\rightarrow$ 38.7\% ($-$0.020 [$-$0.051, +0.007]); the lower bound misses the $-$0.03 margin, so no goal cost is NOT shown. Reviewer workload: 88 held actions (0.20 per episode, 6.2\% of scored actions), 31 of them violations; 0 violations after a block. Exploratory ([released file]): episodes with a held-and-approved action kept their goals (60.8\% vs 58.8\%); the gap is in never-held episodes (38.9\% vs 41.0\%, n = 368), where both arms ran the same way, which we suggest is run-to-run variation against a non-concurrent control (a hypothesis that ``Hold for review'' does not test). Spend \$0.81. \textbf{Deviation 1 (infrastructure, 16:39):} the first start ran 0 episodes. The monitoring server loaded the reader outside its lock, and 6 concurrent first requests raced on the lazy transformers imports (ImportError). Fix: the load moved inside the lock ([released file]; scoring unchanged; regression test [released file]). Rerun from scratch on a GPU job; spend of the failed start \$0.004\\[2pt]
\textbf{tau2-bench, Claude Haiku 5.5}\newline a second outside benchmark: Claude Haiku 5.5 on $\tau$$^2$-bench airline, monitors trained on four frontier agents' labelled trajectories ((i) within-tool reader vs tool identity, (ii) reader vs TF-IDF) & \textbf{not tested}: 9 violations < the preregistered minimum of 15 (Claude violated in 1.5\% of episodes; the 9 come from 3 tasks). Descriptive only: within-tool AUROC reader 0.756, TF-IDF 0.963; pooled reader 0.835, TF-IDF 0.960, tool 0.686; the CAPE-trained probe zero-shot 0.357 pooled ([released file])\\[2pt]
\textbf{Browser-size readers}\newline readers small enough for a browser or a CPU: do they transfer? ((i) browser reader vs TF-IDF on GPT AgentDojo; (ii) non-inferior to the 8B reader on closed CAPE) & \textbf{(i) passes, (ii) fails.} Selected browser reader (best on ``AgentDojo, no labels''): Qwen3-Embedding-0.6B. On ``AgentDojo, closed agent'' (GPT on AgentDojo) it beats TF-IDF: 0.605 vs 0.335 (+0.270 [0.134, 0.451]); on the closed agents' CAPE-Bench actions it is below the 8B reader: $-$0.097 [$-$0.115, $-$0.079] (Claude 0.912, GPT 0.884 vs 0.989, 0.996). Decoder readers on ``AgentDojo, closed agent'': 0.6B 0.651, 1.7B 0.771, 4B 0.696, 8B 0.840, 14B 0.689; none is within 0.02 of the 8B on both sets, so no CPU option is offered. The shipped 8B reproduces its published AUROCs exactly ([released file]). \textbf{Deviation 1:} the analysis crashed after computing, when writing its JSON (a numpy bool); fix: one bool() cast, no computation changed\\[2pt]
\textbf{tau2-bench, two closed agents}\newline $\tau$$^2$-bench with two closed agents (Claude Haiku 5.5, GPT-6-luna) and a persistent customer; ``tau2-bench, Claude Haiku 5.5'''s monitors ((i) reader vs tool within tool, (ii) TF-IDF vs tool) & \textbf{(ii) passes, (i) fails.} Within-tool AUROC: TF-IDF 0.987 (+0.487 [0.428, 0.500] over tool identity), reader 0.738 (+0.238 [$-$0.053, 0.431]); reported: reader $-$ TF-IDF $-$0.250 [$-$0.518, $-$0.077]. Per agent: Claude reader 0.555 / TF-IDF 0.966; GPT 0.739 / 1.000. Reading fixed in advance: on $\tau$$^2$-bench, labels of the policy from other agents transfer to both closed agents through word counts, and the reader does not use them as well ([released file]). Spend (from the ledgers): Claude \$1.71, GPT \$0.83\\[2pt]
\textbf{tau2-bench, unseen tasks}\newline $\tau$$^2$-bench, policy or tasks? ``tau2-bench, two closed agents'''s test, training without the test tasks ((i) reader vs TF-IDF, (ii) reader vs tool; within tool) & \textbf{(i) and (ii) fail: neither monitor transfers to unseen tasks.} Within-tool AUROC: reader 0.535, TF-IDF 0.519 (``tau2-bench, two closed agents'' with the same tasks: 0.738, 0.987); reader $-$ TF-IDF +0.017 [$-$0.514, 0.539]; reader $-$ tool +0.035 [$-$0.230, 0.259]. TF-IDF with entities masked 0.366 (exploratory). ``tau2-bench, two closed agents'''s word-count advantage was memorisation of the tasks ([released file])\\[2pt]
\textbf{tau2-bench, reasoning judge}\newline $\tau$$^2$-bench, a zero-shot reasoning judge (Qwen3-32B, thinking; no labels) on ``tau2-bench, two closed agents'''s writes ((i) judge vs tool, (ii) judge vs ``tau2-bench, unseen tasks'''s reader; within tool) & \textbf{(i) and (ii) fail (intervals include 0).} Within-tool AUROC 0.804 (``tau2-bench, unseen tasks'''s reader 0.535); judge $-$ tool +0.304 [$-$0.084, 0.427]; judge $-$ reader +0.268 [$-$0.074, 0.511]. Pooled 0.785; 53\% caught at the 5\% cut as first computed; 34.5\% with exactly 5\% reviewed (ties at the cut counted in proportion; [released file]); at 50: 12 of 19 caught, 16 of 243 safe flagged. 262/262 answers parsed, none truncated (median 485 tokens). 88\% of scores are exactly 0. ``tau2-bench, Claude Haiku 5.5'': within tool 0.694, 2 of 9 caught at 50. Large point estimate, too few violating tasks to confirm ([released file])\\[2pt]
\textbf{Reasoning judge, frontier agents}\newline $\tau$$^2$-bench, ``tau2-bench, reasoning judge'''s judge unchanged on the 894 published writes of four closed frontier agents ((i) judge vs tool, (ii) judge vs reader cross-validated by task; within tool) & \textbf{(i) passes, (ii) fails.} Within tool: judge 0.766, reader (cross-validated) 0.646, TF-IDF (cross-validated) 0.490; judge $-$ tool +0.266 [0.151, 0.338]; judge $-$ reader (cross-validated) +0.120 [$-$0.047, 0.274]. Reported: 4-sample mean 0.838, $-$ reader (cross-validated) +0.192 [0.022, 0.336]; per agent (judge) 0.768 / 0.771 / 0.775 / 0.618 (o4-mini). With exactly 5\% of writes reviewed: reader (cross-validated) catches 18.4\%, the judge 11.7\%, the 4-sample mean 15.3\%, the tool 13.0\% (most possible 19.7\%). 894/894 parsed in every sample; none truncated; 63\% of scores exactly 0. Reading (fixed): the judge checks rules beyond the tool identity, but not better than the reader trained on other tasks ([released file])\\[2pt]
\textbf{Monitor-aware agent}\newline monitor-aware closed agent: Claude Haiku 5.5 with and without CAPE's oversight notice, 990 scenarios $\times$ 2 arms, no gate; the shipped reader with a cut frozen on Claude's earlier unlabelled traffic ((i) caught aware $-$ plain; (ii) aware caught) & \textbf{(i) and (ii) pass.} Caught at the frozen cut: plain 83.6\%, aware 97.1\% [92.9\%, 100\%]; aware $-$ plain +0.136 [+0.050, +0.225]. Being told did not change how often Claude violated ($-$0.003 [$-$0.012, +0.005]) or reached the goal (37.6\% vs 37.8\%). Note: the notice moves the score distribution up: the frozen cut escalates 15.2\% of the aware arm's actions (plain 5.3\%). At each arm's own 5\% budget: 88.6\% vs 83.6\%. AUROC within template 0.994 vs 0.967 ([released file])\\[2pt]
\textbf{Adaptive attacker}\newline adaptive attacker with score access: Qwen3-32B rewrites the free text of 134 real closed-agent violations (2 rounds $\times$ 16, score feedback), the checker confirms each still violates ((i) reader caught after attack; (ii) reader $-$ TF-IDF) & \textbf{(i) and (ii) fail.} Reader caught 96.3\% $\rightarrow$ 29.1\% [21.6\%, 36.6\%] after attack (Claude 11.8\%, GPT 47.0\%); TF-IDF 57.5\% $\rightarrow$ 34.3\%; reader $-$ TF-IDF $-$0.052 [$-$0.149, +0.037]. Every rewrite still violated (replay checker). The rewrites keep the decisive field (recipient, service and host) and add self-reported approvals ('approved under routine maintenance protocol \ldots{}', '--reason \ldots{}'): the learned monitors read the agent's own claims as evidence. Reading as fixed: the monitor is not robust to adaptive rewording ([released file])\\[2pt]
\textbf{Canonical-action monitor}\newline canonical-action monitor: the shipped reader and recipe trained on canonical actions (free-text fields $\rightarrow$ "[free text]", quoted prose and \# comments removed); ``Adaptive attacker'''s attack re-run against it with seed 1 ((i) caught after attack; (ii) clean catch $-$ shipped on Claude $\cup$ GPT; (iii) AgentDojo within-suite $-$ shipped) & \textbf{(i), (ii) and (iii) fail.} After attack 50.8\% [42.5\%, 59.0\%] caught (shipped reader 29.1\%); before attack 55.2\%. Clean catch at 5\%: Claude 94.1\% $\rightarrow$ 51.5\%, GPT 98.5\% $\rightarrow$ 59.1\%; (ii) $-$0.410 [$-$0.493, $-$0.311]. AgentDojo (iii) $-$0.029 [$-$0.080, +0.010]. By tool: send\_message 29.3\% before and after (the canonical score cannot move, as Corollary 1 predicts); execute\_mock\_command 96.2\% $\rightarrow$ 84.6\% (commands keep their words). Cause: the premise that the label depends on the action only through $\kappa$(A) is false for 81 of the 82 send\_message targets (one family, tp\_campaign\_test: the body carries the draft that is leaked), so the body is payload, not only self-report [this cause was wrong: see "Correction to the ``Canonical-action monitor'' cause" below]. Reading as fixed: canonical mode is an option with its measured cost, not the default ([released file])\\[2pt]
\textbf{Training against the attack}\newline adversarial training: Gemma training actions of the attacked tools, violating AND safe (300 + 300), rewritten by ``Adaptive attacker'''s attacker (8 each), labelled by replay (4,554 rewrites; 28 changed label) and added to training; ``Adaptive attacker'''s attack re-run against the hardened monitor with seed 2 and told the defense ((i) caught after attack; (ii) clean catch $-$ shipped on Claude $\cup$ GPT; (iii) AgentDojo $-$ shipped) & \textbf{(i), (ii) and (iii) fail.} After the defense-aware attack 59.7\% [51.5\%, 67.9\%] caught (before 92.5\%); by tool: send\_message 98.8\% $\rightarrow$ 47.6\%, execute\_mock\_command 82.7\% $\rightarrow$ 78.8\%. The shipped probe on the same rewrites: 96.3\% $\rightarrow$ 57.5\%, so hardening added about 2 points. Clean catch at 5\%: Claude 94.1\% $\rightarrow$ 88.2\%, GPT 98.5\% $\rightarrow$ 97.0\%; (ii) $-$0.037 [$-$0.078, +0.016]. AgentDojo (iii) +0.012 [$-$0.050, +0.071]: the lower bound equals the margin after the locked code's 4-decimal rounding, so it fails as registered. Reading as fixed: augmentation with a known attack does not stop a defense-aware attacker; three preregistered defenses against score-querying rewording have failed ([released file])\\[2pt]
\textbf{Cut set again after a prompt change}\newline set the cut again after a prompt change, live: Claude Haiku 5.5 with ``Monitor-aware agent'''s oversight notice, 400 scenarios $\times$ (ungated, gated); the cut set from ``Monitor-aware agent'''s aware-arm reader scores (3,046, no labels), budget 5\%, frozen ((i) violation rate gated $-$ ungated; (ii) share of ungated actions escalated within [0.025, 0.075]) & \textbf{(i) and (ii) pass.} Episodes with an executed violation 6.5\% (26/400) $\rightarrow$ 0.0\% (0/400): $-$0.065 [$-$0.090, $-$0.043]. The new cut escalated 4.73\% [3.49\%, 6.10\%] of the ungated arm's 1,206 actions; the cut frozen on plain traffic would have escalated 16.6\%. 25 of 54 blocked actions would have been violations; goals 38.5\% $\rightarrow$ 35.5\%; violations caught at the new cut in the ungated arm 88.5\%; median round trip 83 ms. Reading as fixed: setting the cut again on unlabelled traffic restores the budget, and the gate prevents the violations of an agent that knows it is monitored ([released file])\\[2pt]
\textbf{Observable geometry (G1--G3)}\newline training-free observable geometry: on 8 open models ((i) answers vs text, (ii) 8 samples keep $\geq$95\%, (iii) 88 vs 8 questions); on Claude Haiku 5.5 ((i) split-half > 0.7, (ii) vs text); exact Gram recovery on Qwen3-0.6B ((i)--(iii)) & \textbf{all pass}: (i) 8/8 (gains +0.16 to +0.49, lower bounds $\geq$ 0.141); (ii) 8/8 (lower bounds 0.957--0.997); (iii) 8/8; (i) 0.950 [0.939, 0.960]; (ii) +0.438 [0.391, 0.485]; (i) Pearson 1.0 (8 digits) at 2$\times$ d(d+1)/2, median cosine error 7e-6; (ii) error 0.51 $\rightarrow$ 7e-6 across the threshold; (iii) vs baselines 0.40, 0.11. Before the lock, from dev and the pilot only: (ii) changed from an absolute margin to a ratio, (i)'s bound raised 0.5 $\rightarrow$ 0.7 (both disclosed in the prereg). After the lock: Correction 1 to Theorem 5 (below)\\[2pt]
\textbf{Answers vs replies (G5)}\newline which outputs carry the state: yes/no answers vs the model's own replies ((i) exact answers vs 8 replies; (ii) 88 one-token sampled answers vs 1 reply; (iii) replies saturate) & \textbf{all pass}: (i) 8/8 (+0.16 to +0.48); (ii) 7/8 (gpt-oss $-$0.048 [$-$0.075, $-$0.023]); (iii) 8/8 ($\leq$ 0.031). Replies' geometry 0.22--0.44 (0.22--0.34 for seven models; Gemma-3 4B 0.44), close to text. [Corrected 2026-10-22: the row said 0.22--0.34; [released file] has Gemma-3 4B at 0.442.] Dev dry run of the dropped (iv) design disclosed in the prereg.\\[2pt]
\textbf{Shared geometry (G6)}\newline the observable geometry is shared across models ((i) no own-model excess, equivalence $\pm$0.02; (ii) an open model's hidden geometry is the better stand-in; (iii) answers beat text on four new kinds) & \textbf{all pass}: (i) +0.002 [$-$0.001, +0.005]; (ii) 6/7 (Mistral $-$0.049); (iii) 6/8 (Gemma, gpt-oss include 0). On fresh contexts agreement is lower than on ``Interrogation state, open chat models'' (answers 0.22--0.64); Gemma's replies beat its answers (0.43 vs 0.34). Motivated by an independent review's control (recorded in the ledger).\\[2pt]
\textbf{What each API channel reveals (G7)}\newline what each API output channel reveals of a model's own state ((i) answers no excess; (ii) top-20 leaks; (iii) ladder; (iv) log-probabilities leak most; (v) log-probabilities identify the model) & \textbf{all pass}: answers +0.001 (chosen layer), $-$0.006 (last layer); top-1/5/20 +0.056/+0.085/+0.115; full +0.112; log-probability +0.221 [0.211, 0.230]; own model best 7/8 with log-probabilities vs 1/8 with answers.\\[2pt]
\textbf{Reading the reasoning (G8)}\newline does reading an agent's reasoning help a trained monitor? ((i) no in-agent gain on action labels; (ii) the reasoning hurts transfer; (iii) each agent; (iv) later violations) & \textbf{2 of 4 pass.} (i) FAIL: pooled in-agent cot $-$ act $-$0.005 [$-$0.018, +0.013] (margin +0.01); Qwen3.8 gains +0.028 [0.002, 0.078] (73 positives), gpt-oss $-$0.034, Muse $-$0.009. (ii) pass: transfer from Gemma-4 $-$0.033 [$-$0.046, $-$0.017]. (iii) pass, narrowly: $-$0.055, $-$0.0004, $-$0.042. (iv) inconclusive: later violations +0.025 [$-$0.033, +0.085]; AUROC 0.59--0.70.\\[2pt]
\textbf{Live self-policing (G9)}\newline can an agent police itself in real time through its own API? ((i) gate prevents violations; (ii) resampling prevents them with few reviews; (iii) no loss of task success) & \textbf{2 of 3 pass.} Violating episodes: plain 27.9\%; gate@1.0 4.95\% (relative reduction 0.82 [0.76, 0.88]; (i) pass); resampling@0.75 4.35\% (0.84 [0.78, 0.90]) but 0.404 reviews per episode against the 0.40 limit ((ii) FAIL). Success: plain 33.5\%; gate 43.7\% (+10.2 [+6.6, +14.3]); resampling 38.7\% (+5.3 [+1.8, +9.2]) ((iii) pass). Live flag rates on candidate draws: resampling tau 0.94, phi 0.17; gate tau 0.85, phi 0.06.\\[2pt]
\textbf{A closed agent polices itself (G10)}\newline can a closed agent police itself in real time through its own API? ((i) prevents violations its transcript decides; (ii) cannot prevent those decided by an unchecked fact, as Corollary 1 predicts; (iii) no loss of task success) & \textbf{all pass.} Decidable (470 scenarios): violating episodes 2.77\% $\rightarrow$ 0.85\%, reduction 0.69 [0.51, 0.88] ((i)). Unchecked-fact (21 mailing-list scenarios): 100\% $\rightarrow$ 92.9\%, reduction 0.07 [0.00, 0.17] ((ii)). Success 40.4\% $\rightarrow$ 40.1\%, $-$0.3 [$-$2.1, +1.6] ((iii)). All scenarios: 6.9\% $\rightarrow$ 4.8\% (0.31 [0.19, 0.43]); 0.31 reviews per episode.\\[2pt]
\textbf{The prevention law across agents (G11, stopped)}\newline does an agent's offline self-audit predict its live self-policing? (v4fresh) & \textbf{STOPPED} (deviation): after 300 Qwen3.8 episodes neither arm had a violation, so the predicted reduction was untestable; base rates had not been checked; replaced by G11b.\\[2pt]
\textbf{The prevention law across agents (G11b)}\newline the same, on templates where each agent violates ((i) law within 0.15; (ii) Muse > Qwen; (iii) success) & \textbf{1 of 3 pass.} (i) FAIL: Muse-Glimmer 0.23 [0.13, 0.35] vs 0.722; Qwen3.8 $-$0.06 [$-$0.23, 0.09] vs 0.167. (ii) pass: +0.29 [0.10, 0.50]. (iii) FAIL: Qwen3.8 success $-$0.016 [$-$0.048, +0.016]. A note published before any Muse outcome ([a public commit]) traced the Qwen miss to the template mix of the offline sample and gave composition-weighted predictions, 0.273 (Muse, if it does not flag mailing-list sends) and 0.052 (Qwen): errors 0.04 and 0.11 (not a registered endpoint).\\[2pt]
\textbf{The composition-weighted law (G12)}\newline is live self-policing predicted by the composition-weighted law? ((i) each agent within 0.15; (ii) mean error <= 0.10) & \textbf{PASS ((i) 5/5, (ii))}: observed vs predicted reduction: Gemma-4 0.634 vs 0.671; gpt-oss 0.171 vs 0.146; Qwen3.8 0.014 vs 0.033; Muse-Glimmer 0.269 vs 0.254; Claude 0.302 vs 0.381 (difference $-$0.079 [$-$0.128, $-$0.032]: within 0.15 but a systematic over-prediction); mean absolute error 0.035 ([released file])\\[2pt]
\textbf{Self-policing against prompt injection (G13)}\newline does self-policing stop prompt injection live on AgentDojo, by the predicted amount? ((i) Gemma-4 reduction lower >= 0.50; (ii) law within 0.15, both agents; (iii) utility lower > -0.03) & \textbf{4 of 5 pass.} (i) pass: Gemma-4 attack success 11.5\% -> 1.8\%, reduction 0.843 [0.756, 0.912]. (ii): gpt-oss pass (observed 0.502 [0.420, 0.577] vs predicted 0.584); Gemma-4 FAIL (0.843 vs predicted 1.00; miss 0.157 > 0.15). (iii) pass both (utility +0.023 [+0.002, +0.043]; +0.002 [$-$0.027, +0.032]). Exploratory: in 17 of Gemma-4's 24 remaining successes the agent's audit had flagged a call that the stand-in reviewer (injection-only values) let run; 7 (travel) were unflagged, and travel's tau = 1.0 had been assumed in the lock without offline units ([released file])\\[2pt]
\textbf{Recreating the decision state, trained comparison (G14)}\newline can a black-box agent's decision state be recreated through its API well enough to monitor it with a white-box probe? ((i) grey $-$ white > $-$0.03, four open agents; (ii) other agents' coordinates $-$ own white > $-$0.07; (iii) Claude restart AUROC >= 0.80, lower >= 0.65) & \textbf{PASS ((i), (ii), (iii))} --- a TRAINED comparison, not the method: recreation (reader + answers, maps and probe fit on open agents) $-$ true state +0.063 [0.026, 0.122]; other agents' coordinates $-$ own white box +0.009 [$-$0.050, 0.097]; Claude restarts 1.000 (reader alone also 1.000: in this template allowed calls are mostly help/status commands, so the call text separates them). The agent's answers add nothing to the reader ($-$0.000 [$-$0.017, 0.012]) ([released file])\\[2pt]
\textbf{The agent's own answers, training-free (G15)}\newline does a training-free, API-only read-out of an agent's own answers monitor its decisions as well as a probe trained on its hidden states? ((i) mean of 8 $-$ single answer > 0; (ii) mean of 8 $-$ trained white box > $-$0.05; (iii) Claude restart AUROC >= 0.80, lower >= 0.65) & \textbf{PASS ((i), (ii), (iii))}: within-template AUROC, 8 answers (fixed) / 1 answer / probe trained on true states: Gemma-4 0.994/0.966/0.980; gpt-oss 0.822/0.758/0.772; Qwen3.8 0.951/0.894/0.826; Muse-Glimmer 0.929/0.935/0.941. (i) +0.036 [0.006, 0.064]; (ii) +0.045 [$-$0.017, 0.135]; (iii) Claude restarts 0.909 [0.858, 0.952] (single answer 0.323). Claude pooled 0.475: mailing-list sends (fact not in the transcript) mostly not flagged (4.8\% above the median allowed call) ([released file])\\[2pt]
\textbf{Own answers or another model's (G16)}\newline does the monitor need the agent's OWN answers, or does another model's answers to the same questions do as well? ((i) two-sided: own $-$ other, four open agents; (ii) Claude restarts, reported) & \textbf{No detectable difference ((i) includes 0)}: within-template AUROC, own vs another model's answers to the same 8 questions: Gemma-4 0.994 vs 0.993 (Muse-Glimmer answering); gpt-oss 0.822 vs 0.799, Qwen3.8 0.951 vs 0.889, Muse-Glimmer 0.929 vs 0.968 (Gemma-4 answering); mean own $-$ other +0.012 [$-$0.018, 0.040]. (ii) (reported): Claude restarts, own 0.909 vs Gemma-4 answering 1.000, $-$0.091 [$-$0.142, $-$0.045]. Reading fixed in advance: the monitor is a fixed questionnaire that any capable model can answer about the call; the title drops "own" ([released file])\\[2pt]
\textbf{Training-free read-out, prompt injection (G17)}\newline does the training-free read-out of open agents' own answers detect prompt injection on a second benchmark? ((i) within-suite AUROC lower 95\% > 0.60, Gemma-4 and gpt-oss-20b; (ii) read-out $-$ single answer, reported) & \textbf{PASS ((i), both agents)}: within-suite AUROC of the fixed read-out, Gemma-4 0.985 [0.964, 0.998] (816 calls, 216 injected), gpt-oss-20b 0.904 [0.872, 0.932] (821, 221); single answer 0.975 and 0.900; (ii) read-out $-$ single +0.007 [$-$0.004, 0.019]; parsed 94.5\% and 91.1\%. Disclosed: 17 gpt-oss requests failed in the server (HarmonyError) after 3 retries, leaving 39 of 6,568 question slots unanswered (36 on allowed calls; an unanswered question takes the neutral angle, which raises allowed calls' scores, against the read-out); Gemma-4's server died at shutdown, after its last unit ([released file])\\[2pt]
\textbf{Training-free read-out, a second benchmark for Claude Haiku 5.5 (G18)}\newline does the training-free read-out of a CLOSED agent's own answers detect its policy violations on a second benchmark? ((i) within-tool AUROC lower 95\% > 0.50; (ii) read-out $-$ single answer, (iii) read-out $-$ reasoning judge, reported) & \textbf{PASS ((i))}: within-tool AUROC of the fixed read-out 0.704 [0.536, 0.822] (360 writes, 50 violations, 7 tasks with a violation; below the 0.75 "strong" reading); single answer 0.555; (ii) +0.150 [$-$0.004, 0.331]; (iii) vs the Qwen3-32B reasoning judge on its 224 writes 0.742 vs 0.578, +0.164 [$-$0.097, 0.429]; caught at a 5\% budget 20\% vs 4.9\%; parsed 93\%. Exploratory: within task and tool 0.61 [0.46, 0.97] ([released file], [released file])\\[2pt]
\textbf{Own conversation or outside reviewer, same closed model (G19)}\newline does a closed agent's answer need its own conversation? ((i) two-sided: own $-$ outside reviewer, same model, read-out of 8 answers, within tool) & \textbf{No detectable difference ((i) includes 0)}: within-tool AUROC, read-out of 8 answers, own conversation 0.699 vs outside reviewer 0.656, own $-$ outside +0.043 [$-$0.136, 0.226]; single overall question 0.552 vs 0.637, $-$0.086 [$-$0.282, 0.141]; outside answers parsed 99.9\%. Reading fixed in advance: the paper must not claim that "own" matters; the value is in the questions and the read-out ([released file]). Spend \$0.74\\[2pt]
\textbf{Questionnaire or reasoning judge (G20)}\newline does the fixed questionnaire do better than a holistic reasoning judge? ((i) two-sided: read-out $-$ Gemma-4 reasoning judge, mean over four open agents, within template) & \textbf{Questionnaire better ((i) lower > 0)}: within-template AUROC, read-out of the agent's own 8 answers vs Gemma-4 reasoning judge (4 samples): Gemma-4 0.994 vs 0.989, gpt-oss 0.822 vs 0.755, Qwen3.8 0.951 vs 0.780, Muse-Glimmer 0.929 vs 0.880; mean +0.073 [0.021, 0.123]; judge parsed 97--99\%. Reported: Claude restarts, judge 1.000 vs Claude's read-out 0.909 ($-$0.091 [$-$0.143, $-$0.044]); Claude pooled, judge 0.693 vs 0.475 (the judge reasons; Claude's answers had reasoning off) ([released file])\\[2pt]
\textbf{An open model answers about a closed agent (G21)}\newline can an open model answer the questionnaire about a closed agent's calls on the second benchmark? ((i) Gemma-4's read-out about Claude's tau2 writes, within-tool AUROC lower 95\% > 0.50; (ii) own $-$ Gemma-4, reported) & \textbf{PASS ((i))}: Gemma-4's read-out about Claude's tau2 writes 0.989 [0.954, 1.000] within tool (parsed 99.7\%); Claude's own 0.704; (ii) own $-$ Gemma-4 $-$0.285 [$-$0.424, $-$0.171]: the open answerer does better. Not separable here: Gemma-4 answered with reasoning on, Claude (G18, G15's settings) with reasoning off. Exploratory controls: within task and tool 0.965 [0.927, 1.000] (Claude's own 0.61); without the targeted run 0.981; the overall question alone 0.960 ([released file])\\[2pt]
\textbf{One universal questionnaire on two benchmarks (G22)}\newline does one universal questionnaire, written without either benchmark, work on both? (U1a CAPE-Bench, (i') tau2-bench: universal read-out lower 95\% > 0.50; (ii) universal $-$ tailored, same answerer, reported) & \textbf{PASS (U1a, (i'))}: one universal questionnaire, answered by Gemma-4 about the same calls: CAPE-Bench mean over three agents 0.890 [0.818, 0.950] (gpt-oss 0.808, Qwen3.8 0.906, Muse-Glimmer 0.955); tau2-bench 0.935 [0.841, 0.985]. Universal $-$ tailored (same answerer): CAPE-Bench +0.004 [$-$0.029, 0.036] (no detectable difference); tau2-bench $-$0.054 [$-$0.140, $-$0.014] (the tailored questions do better). Reported: Claude restarts 0.948 vs 1.000, pooled 0.443 vs 0.630; parsed 97--98\%. A comparison setting (another model answers), not the method ([released file])\\[2pt]
\textbf{The theory predicts which questions monitor (G23)}\newline does the theory predict, without labels, which questions monitor? ((i) Spearman between each question's label-free Fisher-Rao signal-to-noise ratio and its AUROC, mean over five targets, lower 95\% > 0; (ii), (iii) SNR-weighted read-out vs equal and raw-variance weights, two-sided) & \textbf{(i) PASS; (ii), (iii) no detectable difference}: mean Spearman between each question's label-free Fisher-Rao SNR and its AUROC 0.745 [0.608, 0.779] over five targets (16 questions, Gemma-4 answering); raw-P(yes) variance does as well (0.747; difference $-$0.002 [$-$0.026, 0.015]); SNR-weighted $-$ equal +0.008 [$-$0.177, 0.032]; $-$ raw-variance weights +0.005 [$-$0.177, 0.023]. A comparison setting; G25 is the same test on the agents' own answers ([released file])\\[2pt]
\textbf{The closed agent's own answers with reasoning (G24)}\newline do the closed agent's own answers monitor it when it reasons? ((i) reasoning-on read-out, within-tool AUROC lower 95\% > 0.50; (ii) on $-$ off, two-sided) & \textbf{(i) FAILS; (ii) no detectable difference}: reasoning on 0.568 [0.454, 0.771], reasoning off (same questions, G18's 4 answers each) 0.619; on $-$ off $-$0.051 [$-$0.246, 0.150]; only 70\% of answers parsed (with reasoning, Claude often hedged or wrote to the customer instead of answering; exploratory: the unparsed share does not track violations, AUROC 0.52). Reading fixed in advance: reasoning does not help detectably here; Claude's own answers remain a weak monitor on tau2-bench. Spend \$0.27 ([released file])\\[2pt]
\textbf{The theory predicts which of the agent's own questions monitor (G25)}\newline does the theory predict which questions monitor, on the method's own answers? (G23's (i)-(iii) on each agent's own 8 answers: four open agents and Claude on CAPE-Bench, two agents on AgentDojo, Claude on tau2) & \textbf{(i) PASS; (ii), (iii) opposite}: mean Spearman between each question's label-free Fisher-Rao SNR and its AUROC 0.580 [0.470, 0.648] over eight targets (per target 0.98, 0.93, 0.95, 0.74, $-$0.03, 0.31, 0.40, 0.36); raw-P(yes) variance does as well (0.586; difference $-$0.006 [$-$0.051, 0.036]). SNR weighting lowers the read-out: $-$ equal weights $-$0.068 [$-$0.142, $-$0.047]; $-$ raw-variance weights $-$0.028 [$-$0.102, $-$0.011]. Reading: a label-free measure tells which of the agent's own questions carry the signal, but the noise model adds nothing beyond plain variance here, and weighting by it hurts; equal weights stay ([released file])\\[2pt]
\textbf{The closed model reviews its own calls, with reasoning (G26)}\newline does a closed model, through its own API only, monitor its own calls when it reviews them with reasoning? ((i) within-tool AUROC lower 95\% > 0.50; (ii) $-$ own in-conversation read-out; (iii) $-$ reviewer without reasoning; two-sided) & \textbf{PASS ((i))}: 0.918 [0.730, 0.990] within tool (1.000 within task and tool, reported); flagged 46 of 50 violations and 10 of 130 allowed writes; parsed 100\%; (iii) + 0.262 [0.007, 0.478] over the reviewer without reasoning (0.656); (ii) +0.219 [$-$0.020, 0.427] over Claude's in-conversation read-out (0.699). No other model. Spend \$0.22 ([released file])\\[2pt]
\textbf{State tomography and forecasting (G27)}\newline state tomography and forecasting: does the agent's state, estimated from its own answers at earlier steps (Kalman filter, IV system identification, no labels), forecast its next call's violation before the call is proposed? ((i) forecast $-$ prior, lower 95\% > 0; (ii) filtered $-$ raw; (iii) calibration) & \textbf{(i) FAILS; (ii) no detectable difference; (iii) model misspecified}: on calls after the first step, the forecast from earlier steps ranks the next call at 0.507 against the step-only prior's 0.494 (+0.013 [$-$0.115, 0.155]); filtered $-$ raw $-$0.022 [$-$0.073, 0.030]; NIS/K 0.52 and negative IV-implied noise on 5 of 8 questions: the identified linear dynamics do not fit (they are unstable). Reported: raw pre-call read-out 0.762 and filtered 0.740 within template and tool, but the step index alone scores 0.860 there (G27's episodes end at the first violation). Exploratory, within template, tool and step: raw 0.674 [0.567, 0.776], filtered 0.692 [0.587, 0.801], forecast 0.531 [0.415, 0.657]. Reading: the agent's answers before its call carry information about that call beyond position; earlier steps do not forecast it ([released file], [released file])\\[2pt]
\textbf{The estimated state against the true hidden state (G28)}\newline does the state estimated from the agent's answers track its true hidden state? ((i) RSA of the Kalman estimate $-$ raw answers with the true decision-state geometry, lower 95\% > 0; (ii) RSA of the forecast, lower 95\% > 0) & \textbf{PASS ((i), (ii))}: agreement with the true decision-state geometry on 916 steps: raw 0.048, filtered 0.067, forecast 0.116; filtered $-$ raw +0.019 [0.001, 0.038]; forecast 0.116 [0.078, 0.153]. Exploratory position control: step position alone agrees at 0.210; within same-step pairs (41,607) raw 0.158, filtered 0.180, forecast 0.070, so the forecast's agreement is mostly position, while the answers before the call track the true state within a step ([released file])\\[2pt]
\textbf{State tomography replicated on three agents (G27b)}\newline state tomography and forecasting replicated on three more open agents ((i) pooled forecast $-$ prior, lower 95\% > 0; (ii) pooled filtered $-$ raw; per agent: G28's check against the true state) & \textbf{(i) PASS}: pooled over gpt-oss-20b, Qwen3.8 and Muse-Glimmer, the forecast from the agent's answers at earlier steps beats the prior from position alone, on calls after the first step: +0.209 [0.080, 0.317] (G27 had failed this on Gemma-4); (ii) filtered $-$ raw +0.018 [$-$0.021, 0.062]. Battery sizes: 2,042, 750 and 2,384 steps ([released file])\\[2pt]
\textbf{Live prevention by the estimated state (G29)}\newline live prevention by the agent's estimated state, before the call exists ((i) relative reduction in violating episodes, lower 95\% > 0; (ii) success difference, lower > $-$0.03) & \textbf{Stopped before any episode outcome} (4 minutes after the job started; only the cut, 4.796, was printed): G27 showed that the identified linear dynamics are unstable and that the step index alone ranks violations, so a gate on the Kalman estimate without a position control could not test what it was meant to. Replaced by G29b\\[2pt]
\textbf{Live prevention by the state before the call (G29b)}\newline live prevention by the agent's state read before the call exists, against a position gate and a random gate ((i) pre relative reduction, lower 95\% > 0; (iii) pre $-$ position reduction, lower 95\% > 0; (ii) success, lower > $-$0.03) & \textbf{(i), (iii) FAIL; (ii) pass}: on Gemma-4 (300 scenarios), the gate on the pre-call read-out did not reduce violating episodes (16.7\% $\rightarrow$ 17.3\%; relative reduction $-$0.04 [$-$0.30, 0.16]); the position gate reduced them more (13.3\%; pre $-$ position $-$0.24 [$-$0.50, $-$0.03]); success +0.023 [$-$0.007, 0.053]. The pre-call gate reviewed 5.8\% of calls, and 1.85\% of the calls it held were violations, no more than the random gate (2.8\%): at this cut Gemma-4's pre-call read-out (0.67 within a step, exploratory) ranks too weakly to prevent; GPT-6-luna's (0.92, G40) did (G39). The position gate reviewed 32\% of calls, so (iii) compared gates of very unequal cost ([released file])\\[2pt]
\textbf{The state before a call, beyond position, on three agents (G30)}\newline does the state before a call carry its violation beyond position, and track the true state within a step, on three more agents? ((i) pooled within template x tool x step AUROC of the pre-call read-out, lower 95\% > 0.50; (ii) pooled same-step agreement with the true state, lower 95\% > 0) & \textbf{(i) PASS; (ii) FAIL}: pooled within template x tool x step, the pre-call read-out ranks the coming call's violation at 0.669 [0.594, 0.764] (gpt-oss-20b 0.682, 82 violations in mixed strata; Qwen3.8 0.857, 23; Muse-Glimmer 0.469, 93), confirming Gemma-4's exploratory 0.674 on three more agents; same-step agreement with the true decision-state geometry 0.072 [$-$0.004, 0.144] ([released file])\\[2pt]
\textbf{Repeats against dither: is the within-step effect measurement-limited? (G31)}\newline are the within-step effects limited by the measurement? Re-measure with 16 answers per question: repeats (R) against dither over 4 fixed wordings (D) ((i) pooled within template x tool x step AUROC, D16 $-$ (iv), lower 95\% > 0; (ii) D16 $-$ resampling (16 rollouts), lower 95\% > 0; (iii) same-step agreement with the true state, D16 $-$ (iv), lower 95\% > 0) & \textbf{(i) PASS; (ii), (iii) FAIL}: pooled over 4 agents (Gemma-4 289 steps, gpt-oss-20b 136, Qwen3.8 41, Muse-Glimmer 226), within template x tool x step, the pre-call read-out with 16 dithered answers per question ranks the coming violation at 0.734 against 0.651 with G27's 4 answers ((i) +0.082 [0.001, 0.147]); 16 repeats of one wording do as well (0.728; resampling (16 rollouts) $-$ (iv) +0.076 [0.020, 0.123]; (ii) D16 $-$ resampling (16 rollouts) +0.006 [$-$0.064, 0.066]): the within-step effects are limited by the number of answers, and on open agents dither adds nothing over repeats. Same-step agreement with the true state D16 $-$ (iv) +0.053 [$-$0.014, 0.129] ((iii)). Dither lowers unanimity (Gemma-4 0.87 $\rightarrow$ 0.58); split-half reliability of the read-out 0.85--0.99 ([released file])\\[2pt]
\textbf{Early warning and the horizon of the state (G31b)}\newline early warning and horizon: does the state read before the first call warn of a violation that comes later, and does the information arrive at the call's step? ((i) pooled within template x first-tool AUROC of the first-step read-out for 'violates later', lower 95\% > 0.50; (i) pooled AUROC of the read-out before call t minus before call t$-$1, lower 95\% > 0) & \textbf{(i) FAIL; (i) PASS}: before the first call, the answers do not warn of a violation that comes later (pooled early-warning AUROC 0.498 [0.405, 0.585]; Gemma-4 0.453, reported); the answers just before a call rank it better than the answers one step earlier (+0.163 [0.029, 0.302]): the deciding information enters at the call's step ([released file])\\[2pt]
G31c how long does the state read through the API remember, and is the noise G27's filter assumed right? Label-free ((i) pooled lag-1 autocorrelation of the answer angles within template x step, lower 95\% > 0; (ii) lag-1 $-$ lag-3, lower 95\% > 0; (iii) within-cell variance / nominal noise 1/(m + 1/2), upper 95\% < 1) & \textbf{PASS ((i), (ii), (iii))}: the state read through the API has a short memory: pooled within-episode autocorrelation 0.132 [0.090, 0.171] at one step, 0.042 at two, 0.013 at three (decay 0.119 [0.075, 0.155]); the angles' variance is 0.31 [0.29, 0.34] of the nominal noise 1/(m + 1/2) that G27's filter assumed. Together: forecasts must use the last step or two, and a filter must not use the nominal noise. An earlier run of this analysis on gpt-oss-20b alone (a waiter fired before the other batteries existed; [released file]) is kept as a partial file and is not this result ([released file])\\[2pt]
\textbf{A label-free signal says where the read-out works (G32)}\newline does a label-free quantity from the agent's own answers say where the read-out works? ((i) Spearman over (target, stratum) cells between the read-out's SNR beyond sampling noise and its AUROC, lower 95\% > 0; (ii) AUROC of cells with SNR > 0 minus SNR $\leq$ 0, lower 95\% > 0) & \textbf{PASS ((i), (ii))}: Spearman 0.710 [0.436, 0.870]; cells with SNR > 0 (23) average AUROC 0.961, cells with SNR $\leq$ 0 (6) 0.678, difference +0.282 [0.169, 0.406]. Partial Spearman given cell size and base rate 0.677; raw variance alone 0.669; 82.8\% of cells within the noise ceiling + 0.05. Per target, Claude Haiku 5.5 has the lowest SNR (tau2-bench $-$0.48, AUROC 0.705; CAPE-Bench $-$0.18): its 4 answers agree on 96--97\% of questions, so each question gives about one bit ([released file])\\[2pt]
\textbf{The closed agent's answers, dithered (G33)}\newline the closed agent's own answers, dithered: 4 fixed wordings per question on tau2-bench airline ((i) within-tool AUROC D16 $-$ G18's 4-answer read-out, lower 95\% > 0; (ii) D16 $-$ resampling (16 rollouts), lower 95\% > 0) & \textbf{FAIL ((i), (ii))}: dithered wordings changed the measurement as predicted (unanimity 0.958 $\rightarrow$ 0.781; label-free SNR $-$0.48 $\rightarrow$ 0.73) but did not raise the ranking: within-tool AUROC D16 0.720 [0.573, 0.824], G18's 4 answers 0.705, (iv) 0.690, resampling (16 rollouts) 0.686; (i) D16 $-$ G18\_R4 +0.015 [$-$0.089, 0.104]; (ii) D16 $-$ resampling (16 rollouts) +0.034 [$-$0.088, 0.157]. On this benchmark the closed agent's read-out is not limited by the quantized answers; 16 repeats also raised the nominal SNR (0.88) without raising AUROC, so a positive SNR there does not imply a better ranking. Cost \$11.77 ([released file])\\[2pt]
\textbf{Triangulation: the closed model's own hidden geometry (G34)}\newline triangulation: does a closed model's own answer geometry reconstruct its own unseen hidden geometry? ((i) leave-one-open-model-out mean absolute error of the predicted own agreement < 0.10; (ii) observed inside the predicted interval for $\geq$ 6 of 8; (i) the closed model's triangulated own agreement, lower 95\% > 0) & \textbf{PASS ((i), (ii), (i))}: leave-one-out error 0.061 [0.051, 0.070]; 7 of 8 inside the interval (miss: gpt-oss-20b, whose hidden geometry is the least shared, loading 0.58); Claude Haiku 5.5's own answers against its own hidden geometry, triangulated: 0.495 [0.437, 0.554], within the open models' own range (0.48--0.62). Reported: on (vi)'s published matrices (1,000 contexts of 10 kinds, 88 exact answers) the same leave-one-out is less accurate (error 0.114; 4 of 8 inside) ([released file])\\[2pt]
\textbf{Does the monitor come from the reconstructed state? (G35)}\newline does the monitor come from the reconstructed state? ((i) Spearman between the read-out and the hidden state's violation coordinate within template and label, lower 95\% > 0; (ii) answer vs decision-state geometry over same-template, same-label pairs, lower 95\% > 0; (iii) share of the read-out's AUROC lost when the part the hidden coordinate predicts is removed, lower 95\% > 0) & \textbf{(iii) PASS; (i), (ii) FAIL}: 20\% [13\%, 27\%] of the read-out's AUROC above chance runs through the hidden state's violation coordinate; within a template and label the 4-answer read-out does not follow that coordinate ((i) 0.044 [$-$0.008, 0.115]; (ii) $-$0.011 [$-$0.058, 0.027]). Within such a cell the answers are nearly unanimous, so the read-out barely varies ([released file])\\[2pt]
\textbf{With dither, does the read-out follow the hidden state? (G35b)}\newline with a better measurement, does the read-out follow the hidden state's violation coordinate beyond the label? ((i) ``Interrogation state, Othello'' of the dithered read-out D16 within template x tool x step x label, lower 95\% > 0; (ii) ``Interrogation state, Othello''(D16) $-$ ``Interrogation state, Othello''((iv)), lower 95\% > 0) & \textbf{FAIL ((i), (ii))}: even with dithered answers, within template x tool x step x label the read-out does not follow the hidden state's violation coordinate (``Interrogation state, Othello'' of D16 $-$0.077 [$-$0.250, 0.195]; D16 $-$ (iv) $-$0.041 [$-$0.274, 0.201]); cells are small (28--217 steps with a state per agent) ([released file])\\[2pt]
\textbf{Triangulation at the agent's decision (G36)}\newline triangulation at the agent's decision: how well do an agent's own answers at its calls reconstruct its own decision-state geometry? ((i) leave-one-open-agent-out mean absolute error < 0.10; (ii) inside the interval for $\geq$ 3 of 4; (i) Claude Haiku 5.5's triangulated own agreement, lower 95\% > 0) & \textbf{(i), (i) PASS; (ii) FAIL}: at the agent's decision, an agent's own 8 policy answers agree only weakly with its own decision-state geometry over all pairs of calls (observed 0.010, 0.039, 0.054, 0.161 for Gemma-4, gpt-oss-20b, Qwen3.8, Muse-Glimmer); the triangulation predicts these within 0.045 [0.031, 0.065] on average ((i)) but inside its interval for only 1 of 4 ((ii)); Claude Haiku 5.5's estimate 0.103 [0.085, 0.122] ((i)). Eight policy questions measure a narrow, violation-relevant slice of a rich decision state; the read-out uses that slice (G35: 20\% of its power runs through the hidden violation coordinate), while the broad geometry is reconstructed with many generic questions ((i), G34) ([released file])\\[2pt]
\textbf{The closed agent's state before its call (G37)}\newline the closed agent's state before its call exists, on tau2-bench airline ((i) within-tool AUROC of the pre-call read-out of its own answers about its next action, lower 95\% > 0.50; (ii) dither (iv) $-$ (iv) at equal budget, lower 95\% > 0) & \textbf{FAIL ((i), (ii))} (run under the recorded deviation, [released file]: (iv) on all 360 writes, the dither arm on 288 before the fixed credit ran out): before its write exists, Claude Haiku 5.5's answers about its next action do not rank the coming violation (within-tool AUROC 0.536 [0.434, 0.682]; at the call, G18: 0.705; pre $-$ at call $-$0.169 [$-$0.289, $-$0.027]). On the 288 writes with both arms (33 violations), dither 0.652 vs repeats 0.525, (ii) +0.127 [$-$0.045, 0.253]. Unanimity before the call 0.96. Cost about \$1.9 ([released file])\\[2pt]
\textbf{A second closed agent, at its call and before it (G38)}\newline a second closed agent (GPT-6-luna) read through its own API at its call and before it, on CAPE-Bench ((i) within template x tool AUROC of the at-call read-out, lower 95\% > 0.50; (i) of the pre-call read-out, lower 95\% > 0.50) & \textbf{PASS ((i), (i))}: GPT-6-luna's own answers rank its violations within template x tool at the call 0.979 [0.946, 0.998] and before the call exists 0.765 [0.650, 0.909] (560 calls, 66 violations, 216 scenarios; cost \$1.14). Unanimity 0.85 at the call, 0.77 before it; at the call it abstained on about half the rule questions it could not check ("I can't determine..."), which take the neutral angle. Exploratory position control: the call's position alone ranks at 0.763 within template x tool, the same as (i); within template x tool x position (12 strata with both classes, 27 violations) the pre-call read-out ranks at 0.927 [0.828, 0.989] and the at-call read-out at 0.964 ([released file])\\[2pt]
\textbf{Live prevention on a closed agent before the call (G39)}\newline live prevention on a closed agent (GPT-6-luna) by its state read before the call exists, against a position gate ((i) relative reduction of violating episodes, pre vs none, lower 95\% > 0; (iii) pre $-$ position, lower 95\% > 0; (ii) success, lower > $-$0.05; on the transcript-decidable templates) & \textbf{(i) PASS; (iii), (ii) FAIL}: on the transcript-decidable templates (84 scenarios), the gate on GPT-6-luna's pre-call state cut violating episodes from 26.2\% to 9.5\% (relative reduction 0.636 [0.250, 0.875]); the position gate cut them to 15.5\% (0.409 [0.000, 0.688]); pre $-$ position +0.227 [$-$0.190, 0.625] ((iii) fails: not distinguishable from position at this power); success +0.012 [$-$0.083, 0.107] ((ii) fails: the lower bound is below $-$0.05). Reviews per episode: pre 1.36, position 0.92. Unchecked-fact template (tp\_campaign\_test, 42 scenarios): 95.2\% $\rightarrow$ 97.6\%, no reduction, as predicted. All 378 episodes ran; cost \$1.13 ([released file])\\[2pt]
\textbf{The pre-call gate against a matched position gate (G39b)}\newline the closed agent's pre-call gate against a position gate of matched review rate ((iii) relative reduction pre $-$ position, lower 95\% > 0; (i) pre vs none, lower 95\% > 0; (ii) success, lower > $-$0.05) & \textbf{PASS ((iii), (i), (ii))}: at a matched review rate (pre 1.48, position 1.52 reviews per episode), on 84 decidable scenarios x 3 fresh seeds (252 runs per arm, all complete), the gate on GPT-6-luna's pre-call state cut violating episodes from 31.7\% to 7.9\% (relative reduction 0.750 [0.605, 0.863]) and the position gate to 17.9\% (0.438 [0.256, 0.605]); pre $-$ position +0.313 [0.062, 0.543]; success unchanged (18.7\% in both; difference 0.000 [$-$0.048, 0.048]). Cost \$2.94 ([released file])\\[2pt]
\textbf{The relevant directions of the call (G41)}\newline read the closed agent's state along the directions its call can violate ((i) within-tool AUROC of the policy-relevant read-out minus the full read-out, dithered answers, lower 95\% > 0; (ii) the same on G18's answers) & \textbf{FAIL ((i), (ii))}: relevant $-$ full $-$0.014 [$-$0.059, 0.038] (dithered) and +0.008 [$-$0.021, 0.026] (G18's answers); every read-out 0.70--0.72. Diagnosis across G18, G24, G26, G33, G37, G41 (exploratory): answers given from the agent's own conversation stay near 0.70 whatever the measurement (more answers, dither, relevant questions, reasoning on 0.57, before the call 0.54, where the general question is inverted, 0.434), while the same model as an outside reviewer with reasoning reaches 0.918 (G26): the information is in the transcript, and the agent's own-seat answers are committed to its plan ([released file])\\[2pt]
\textbf{The prevention law for pre-call gates (G42)}\newline does the prevention law predict a live pre-call gate on a closed agent? ((i) |G39b's pre reduction $-$ offline share 0.750| $\leq$ 0.15; (ii) |G39b's position reduction $-$ 0.167| $\leq$ 0.15) & \textbf{(i) PASS; (ii) FAIL}: the offline share of violating calls held predicted the pre-call gate's live reduction exactly (predicted 0.750, observed 0.750 [0.605, 0.863]; error 0.000). It under-predicted the position gate (predicted 0.167, observed 0.438 [0.256, 0.605]; error 0.271), as in G39 (0.125 vs 0.409): the law does not carry to a gate on turn index for this agent ([released file])\\[2pt]
\textbf{A broader battery at the decision (G43)}\newline does a broader battery of the agent's own answers (32 situation questions) reconstruct its decision state? ((i) pooled Spearman with the decision-state geometry, 32 questions $-$ 8 policy questions, lower 95\% > 0; (ii) 32 questions, lower 95\% > 0.2) & PENDING\\[2pt]
\textbf{Self against other forecast (G44)}\newline does an agent's own state forecast its next violation better than another model reading the same conversation? ((i) GPT-6-luna, within template|tool, own $-$ other pre-call AUROC, lower 95\% > 0; (ii) three open agents pooled, within template x tool x step, lower 95\% > 0) & \textbf{(i) PASS; (ii) FAIL}: before its call exists, GPT-6-luna's own answers about its next action forecast its next violation better than Gemma-4 answering the same questions in its seat on the same conversations: within template|tool 0.840 vs 0.728, own $-$ other +0.111 [0.041, 0.185] (402 calls, 27 violations); within template|tool|position 0.923 vs 0.847, +0.076 [$-$0.017, 0.230]. Open agents, pooled within template x tool x step, own $-$ other $-$0.030 [$-$0.159, 0.082]: Qwen3.8 own 0.857 vs other 0.643; gpt-oss-20b 0.682 vs 0.841; Muse-Glimmer 0.468 vs 0.613 ([released file])\\[2pt]
\textbf{Is the self advantage specific to forecasting? (G44b)}\newline is the agent's advantage over another reader specific to forecasting? ((i) pooled over three open agents: (own $-$ other) before the call $-$ (own $-$ other) at the call, lower 95\% > 0) & \textbf{FAIL ((i))}: pooled over the three open agents, (own $-$ other) before the call minus at the call $-$0.045 [$-$0.178, 0.082]: Qwen3.8 +0.152 (at the call +0.062, before +0.214), gpt-oss-20b $-$0.182 (+0.023, $-$0.159), Muse-Glimmer $-$0.106 ($-$0.040, $-$0.145). The agent-specific forecasting advantage holds for the closed agent (G44 (i)) and for Qwen3.8, not for gpt-oss-20b or Muse-Glimmer, whose own pre-call answers rank weakly ([released file])\\[2pt]
\textbf{A closed agent's pre-call state beyond position, confirmed (G40)}\newline confirm on fresh episodes that a closed agent's state before its call ranks the coming violation beyond position ((i) within template x tool x position AUROC of the pre-call read-out, lower 95\% > 0.50; (ii) within template x tool, lower 95\% > 0.50) & \textbf{PASS ((i), (ii))}: on 252 fresh episodes (1,047 calls, 68 violations), GPT-6-luna's answers about its next action, read before the call exists, rank the coming violation at 0.923 [0.811, 0.992] within template x tool x position (9 strata with both classes, 27 violations; G38's exploratory value was 0.927) and at 0.840 [0.732, 0.946] within template x tool, where position alone ranks at 0.777. Cost \$2.06 ([released file])\\[2pt]
\bottomrule
\end{longtable}}

\paragraph{Correction to one recorded cause.} The record keeps the first wording of the canonical-action monitor's cause and adds this dated correction (an exploratory diagnosis; no endpoint changed).
\begin{itemize}
\item The ``Canonical-action monitor'' row first said that the message body "carries the leaked draft", so free text is payload. Checking the data showed this is wrong:
\item \textbf{What decides the label.} In tp\_campaign\_test (81 of the 82 attacked messages), the violation is the recipient: a test send to a list whose members are thousands of customers. That fact appears only in a \texttt{list-info} result, and neither closed agent ever called \texttt{list-info}. CAPE's checker labels the send by its recipient, whatever the body says.
\item \textbf{What the attacker did.} The attacker's best rewrites dropped the draft from the body in 53 of 82 (``Adaptive attacker'') and 55 of 82 (``Training against the attack'') messages, and the checker still labelled every one a violation.
\item \textbf{Training coverage.} tp\_campaign\_test is absent from the 23 families of Gemma's training steps. Among open agents, only Qwen has violations of it (7).
\item \textbf{The cause.} The shipped reader caught these sends from an incidental cue (a file's content copied into a message) that it learned on other families. The attacker can remove that cue without changing the label. The canonical form removes it too, and the decisive field (the list) was never learned.
\item \textbf{A theorem (paper, "a ceiling under rewording").} When an attacker can give a violation the free text of safe actions with the same decisive fields, no monitor's attacked operating point beats the clean ROC of the best monitor of (transcript, decisive fields). The canonical reader's clean catch (51.5\% / 59.1\%) is our estimate of that ceiling. ``Training against the attack'''s 59.7\% after the attack and ``Adaptive attacker'''s 29.1\% are consistent with it.
\item \textbf{The headline without this family.} The shipped reader's exact-budget catch at 5\% is 92.3\% (Claude, 26 violations) and 96.3\% (GPT, 27 violations), against 94.1\% and 98.5\% with it.
\end{itemize}


\end{document}
